% --- Core ---
\usepackage[utf8]{inputenc}
\usepackage[T1]{fontenc}
\usepackage{lmodern}
\usepackage[margin=2.5cm]{geometry}
\usepackage{enumitem}

% --- Math & Symbols ---
\usepackage{amsmath, amssymb, amsfonts}

% --- UI & Layout ---
\usepackage{parskip} % Spasi antar paragraf, tanpa indentasi
\usepackage{float}
\usepackage{graphicx}
\usepackage{caption}
\usepackage{subcaption} % Untuk sub-figures
\usepackage{booktabs}   % Tabel cantik ( \toprule, \midrule, \bottomrule )
\usepackage{array}
\usepackage{tabularx}
\usepackage{multirow}
\usepackage{longtable}

% --- Colors & Links ---
\usepackage[dvipsnames]{xcolor}
\usepackage{hyperref}

% --- Code Snippets (Listings) ---
\usepackage{listings}
\definecolor{codegreen}{rgb}{0,0.6,0}
\definecolor{codegray}{rgb}{0.5,0.5,0.5}
\definecolor{codepurple}{rgb}{0.58,0,0.82}
\definecolor{backcolour}{rgb}{0.95,0.95,0.92}

\lstset{
    backgroundcolor=\color{backcolour},
    commentstyle=\color{codegreen},
    keywordstyle=\color{magenta},
    numberstyle=\tiny\color{codegray},
    stringstyle=\color{codepurple},
    basicstyle=\ttfamily\small,
    breakatwhitespace=false,
    breaklines=true,
    captionpos=b,
    keepspaces=true,
    numbers=left,
    numbersep=5pt,
    showspaces=false,
    showstringspaces=false,
    showtabs=false,
    tabsize=4
}

% --- Diagrams (TikZ) ---
\usepackage{tikz}
\usetikzlibrary{arrows.meta, shapes.geometric, shapes.misc, positioning, calc, fit, backgrounds, matrix, decorations.pathreplacing}

\tikzset{
    blok/.style={rectangle, draw, rounded corners=2pt, fill=blue!6, align=center, minimum height=0.8cm, font=\small},
    blokgol/.style={blok, fill=green!12},
    blokspn/.style={blok, fill=orange!14},
    blokkunci/.style={blok, fill=violet!10},
    data/.style={rectangle, draw, fill=gray!8, align=center, minimum height=0.7cm, font=\small},
    keputusan/.style={diamond, draw, aspect=2, fill=yellow!12, align=center, font=\small, inner sep=1pt},
    xorop/.style={circle, draw, inner sep=0pt, minimum size=0.45cm, path picture={
        \draw (path picture bounding box.north) -- (path picture bounding box.south)
              (path picture bounding box.west) -- (path picture bounding box.east);}},
    addop/.style={rectangle, draw, inner sep=0pt, minimum size=0.42cm, path picture={
        \draw (path picture bounding box.north) -- (path picture bounding box.south)
              (path picture bounding box.west) -- (path picture bounding box.east);}},
    panah/.style={-Stealth, thick},
    garis/.style={thick},
}

% --- Header & Footer ---
\usepackage{fancyhdr}
\pagestyle{fancy}
\fancyhf{}
\fancyhead[L]{\small IF4020 Kriptografi -- \algname}
\fancyhead[R]{\small \leftmark}
\fancyfoot[C]{\thepage}
\setlength{\headheight}{14pt}

% --- Referencing ---
\usepackage[noabbrev, capitalize]{cleveref}
\graphicspath{{images/}}

\renewcommand{\contentsname}{Daftar Isi}
\renewcommand{\listfigurename}{Daftar Gambar}
\renewcommand{\listtablename}{Daftar Tabel}
\renewcommand{\chaptername}{Bab}
\renewcommand{\appendixname}{Lampiran}
\renewcommand{\figurename}{Gambar}
\renewcommand{\tablename}{Tabel}
\renewcommand{\bibname}{Daftar Pustaka}
\renewcommand{\lstlistingname}{Kode}

\crefname{chapter}{Bab}{Bab}
\crefname{section}{Subbab}{Subbab}
\crefname{subsection}{Subbab}{Subbab}
\crefname{figure}{Gambar}{Gambar}
\crefname{table}{Tabel}{Tabel}
\crefname{equation}{Persamaan}{Persamaan}
\crefname{lstlisting}{Kode}{Kode}
\crefname{listing}{Kode}{Kode}
\crefname{appendix}{Lampiran}{Lampiran}
\crefname{enumi}{Langkah}{Langkah}

% --- Notasi ---
\newcommand{\algname}{Pregnant\_\allowbreak Sister}
\newcolumntype{Y}{>{\raggedright\arraybackslash}X}
\setlength{\emergencystretch}{3em}
\newcommand{\isi}[1]{\textcolor{red}{\textbf{[#1]}}}
\newcommand{\code}[1]{\texttt{#1}}
\newcommand{\rotl}{\mathbin{\lll}}

% build-remote.bat mengemas src/ di samping main.tex; kompilasi lokal membaca ../../src
\IfFileExists{src/feistel.py}{\def\srcdir{src}}{\def\srcdir{../../src}}
